\section{Conductor conventions and actual order periods}\label{sec:conventions}
\subsection{Four numbers, not one unnamed conductor}
Let $F$ be a nonarchimedean local field with residue cardinality $q$, and let $V$ be a complex Weil representation with finite inertia image and no monodromy. Its ordinary Artin conductor is
\[
 a(V)=\dim V-\dim V^{I_F}+\Sw(V).
 \tag{22.1}\label{eq:artin}
\]
We use $a$ for the ordinary conductor throughout this part. The Swan conductor, an epsilon exponent for a specified additive character, and the self-pair conductor of a representation of $\GL_n(F)$ are named separately.

Take $\psi_0$ trivial on $\OO_F$ and nontrivial on $\pp_F^{-1}$. Set $\psi_1(x)=\psi_0(\varpi^{-1}x)$, so $\psi_1$ is trivial on $\pp_F$ and nontrivial on $\OO_F$. Under the usual change-of-character convention, if the exponent is written as a power of $q^{1/2-s}$, then
\[
 f_{\psi_0}(V)=a(V),\qquad f_{\psi_1}(V)=a(V)-\dim V.
 \tag{22.2}\label{eq:psi}
\]
The determinant factor in the change of additive character is independent of $s$ and does not alter this exponent shift. Equation \eqref{eq:psi} is a convention translation, not an assertion that every epsilon constant is determined by its exponent.

For a wild-irreducible source $\rho$ of dimension nine with Swan conductor $24$, the audited self-pair calculation below gives the following table.
\begin{center}\begin{tabular}{lrrrrr}\toprule
Object & Dimension & $I_F$-fixed & Swan & Ordinary $a$ & $f_{\psi_1}$\\\midrule
$\rho$ & 9 & 0 & 24 & 33 & 24\\
$\End\rho$ & 81 & 1 & 152 & 232 & 151\\
$\Ad\rho$ & 80 & 0 & 152 & 232 & 152\\\bottomrule
\end{tabular}\end{center}
Here $\Ad\rho$ denotes the trace-zero summand. Over the complex coefficient field, $\End\rho=\C\oplus\Ad\rho$. The scalar summand is unramified, so it contributes no ordinary conductor. It does contribute one to the dimension shift in \eqref{eq:psi}.

Thus the number $152$ has two legitimate appearances, as the Swan conductor of $\End\rho$ and as the $\psi_1$ exponent of its trace-zero summand. It is \emph{not} the $\psi_1$ exponent of all of $\End\rho$, which is $151$. The formal-degree parameter $\delta$ appearing later also equals $152$ in this particular torsion-one comparison. That numerical equality does not make the three definitions interchangeable.

\source{The ordinary conductor identity and additive-character normalization underlie [Q14]. [BH99, \S6.1, Lemma 2] uses the character trivial on $\pp_F$. [AP, \S4 and Theorem 6.9] uses the ordinary pair conductor. The table is an explicit reconciliation of these conventions using the source calculation in Section~\ref{sec:ranknine}.}

\subsection{The principal-order grade calculation}
Let $\A\subset M_n(F)$ be a principal hereditary order of period $e$, with radical $\Pp$. Then $e\mid n$. A principal lattice chain has $e$ successive residue pieces, each of dimension $m=n/e$ over $k_F$. The homogeneous grade of an endomorphism is the direct sum of $e$ matrix blocks of size $m$ by $m$. Consequently
\[
 \dim_{k_F}\Pp^j/\Pp^{j+1}=e m^2=\frac{n^2}{e},
 \qquad [\Pp^{-r}:\A]=q^{r n^2/e} \quad(r>0).
 \tag{22.3}\label{eq:grade}
\]
The period here belongs to the actual principal order carrying the stratum. It cannot be borrowed from a more convenient unrelated lattice chain.

For an actual cuspidal representation containing a positive-level simple stratum $[\A,r,0,\beta]$ to which the hypotheses of [BH99, \S6.1] apply, its Lemma 2 gives
\[
 f_{\psi_1}(\pi)=\frac{nr}{e},\qquad
 a(\pi)=n+\frac{nr}{e}.
 \tag{22.4}\label{eq:principal-conductor}
\]
The first number is the native epsilon exponent in that source. In the present wild cuspidal comparison it equals the standard Swan conductor. A lower critical level in a defining sequence is not the leading level $r$ in this formula.

\begin{proposition}[Necessary period divisibility]\label{prop:period}
Suppose an actual type satisfying \eqref{eq:principal-conductor} is to match a source of rank $n$ and Swan conductor $w>0$. Its principal period and leading level necessarily satisfy
\[
 \frac{n}{\gcd(n,w)}\mid e\mid n,
 \qquad r=\frac{ew}{n}.
 \tag{22.5}\label{eq:period}
\]
These are necessary numerical conditions, not sufficient existence conditions for a simple stratum or a prescribed source.
\end{proposition}
\begin{proof}
The principal condition gives $e\mid n$. Matching \eqref{eq:principal-conductor} gives $nr=ew$, so $n\mid ew$. Writing $n=gn'$ and $w=gw'$ with $\gcd(n',w')=1$ yields $n'\mid e$. Conversely these divisibilities make $ew/n$ an integer. They say nothing about purity, intertwining, or character occurrence.
\end{proof}

For $(n,w)=(9,24)$, the possibilities are $(e,r)=(3,8)$ and $(9,24)$. A maximal-order construction of period one cannot have this conductor. The earlier unramified depth-two candidate has Swan conductor $18$ and ordinary conductor $27$, not $24$ and $33$. Replacing the leading level eight by a critical level four in the period-three formula would create a different error rather than repair the comparison.

\subsection{A self-twist check that survives conductor matching}
A second necessary invariant is the torsion number, the order of the group of unramified self twists. For a maximal simple type, the exact formula in [BK, 6.2.5] identifies that number as $n/e(K/F)$ for the actual parameter field $K$.

\begin{lemma}[Wild irreducibility excludes nontrivial unramified self twists]\label{lem:source-twist}
If $\rho|_{P_F}$ is irreducible and an unramified character $\chi$ satisfies $\rho\otimes\chi\simeq\rho$, then $\chi=1$.
\end{lemma}
\begin{proof}
An intertwiner restricts to an endomorphism of the irreducible $P_F$ representation, because $\chi$ is trivial on $P_F$. It is therefore a nonzero scalar. Applying the intertwining identity to every Weil element forces $\chi$ to be one.
\end{proof}

A period-three candidate may pass the first conductor test and still fail this second test if its actual maximal parameter field has residue degree three. That candidate has torsion number three, while the prescribed source has torsion number one. This is why the conductor equation does not complete the arithmetic attachment.
